\section{Finite markings and integral models}
\label{sec:finite-markings}

This section provides the finite algebraic data used by both sides of
the proof. The whole models give coefficients whose boundary layers
belong to the next height stratum; their localized summands give the
marking covers whose completed tower is studied geometrically. We also
identify the full division-point lift schemes, because later transfer
maps require their scheme structure and natural group actions.

Fix a prime $p$ and a height $n$.  Choose a finite field $\kk$ containing
the fields of definition of the Honda groups $\Gamma_j$, $1\leq j\leq n$,
and of their endomorphisms.  We use Honda coordinates for which
$[p]_{\Gamma_j}(T)=T^{p^j}$.  In this section
\[
 G=P_n,\qquad
 A_m=\kk[[u_m,\ldots,u_{n-1}]],\qquad 1\leq m\leq n.
\]
The formal group $\cF$ over $A_m$ is the restriction of the universal
deformation of $\Gamma_n$ to the height-at-least-$m$ locus.
The deformation-ring construction is due to
\citet[Proposition~1.1 and Theorem~3.1]{lubinTate1966}; the Honda conventions and
endomorphism structure are recalled in \citet[Appendix~A2.2]{ravenel1986}.  Its finite
division schemes are formed over $A_m$ before any localization.  If
$m<n$, put
\[
 x=u_m,\qquad B_m=A_m[1/x],\qquad r=n-m.
\]
The chosen parameters are successive Hasse coefficients: the ideals of
the deeper height loci are $(u_m,\ldots,u_{j-1})$, and the first
coefficient of $[p]_{\cF}(T)$ over $A_m$ is $xT^{p^m}$; compare
\citet[(1.1) and \S2]{strauch2007}.  Replacing a
parameter by a unit multiple gives the same constructions.

The action of $G$ comes with coordinate changes
$\alpha_g:g^*\cF\longrightarrow\cF$.  Write $\chi_g=\alpha_g'(0)$ and
let $\cL_m$ denote the free rank-one $A_m$-module with this semilinear
action.  Composition of the $\alpha_g$ gives the cocycle identity for
$\chi_g$.  Comparing the leading coefficients of the $p$-series gives
\begin{equation}
 g(x)=\chi_g^{p^m-1}x.
 \label{eq:mark-height-character}
\end{equation}
Thus $\cL_m$ is the periodicity line, up to the immaterial choice of
inverse convention.  All its integer powers will be retained.  All
finite $A_m$-modules below carry their compact topology; localizations
carry the bounded-denominator coefficient convention of
\cref{sec:continuous}.

\subsection{The coefficient class}

Since the $p$-series factors through $T^{p^m}$, its $\ell$-fold iterate
factors through $T^{p^{m\ell}}$.  For $\ell\geq1$ and
$0\leq b\leq m\ell$, define
\[
 D_m(\ell,b)
   =A_m[\,T^{p^b}\,]\ \subseteq\
     A_m[[T]]/([p^\ell]_{\cF}(T)).
\]
The brackets here mean the $A_m$-subalgebra generated by the indicated
element; it will be finite over $A_m$.

\begin{definition}[Allowed coefficients]
\label[definition]{def:mark-coefficients}
The basic coefficients at height $m$ are
\begin{equation}
 \cL_m^{\otimes w}\otimes_{A_m}
       \bigotimes_{j=1}^{s}D_m(\ell_j,b_j),
 \qquad w\in\Z,\quad 0\leq b_j\leq m\ell_j,
 \label{eq:mark-basic-coefficients}
\end{equation}
where $s$ may be zero.  These are functions on whole tuples of division
points, with the indicated powered coordinates; no independence
condition is imposed.  The allowed class $\mathscr C_m$ is their closure
under finite direct sums, tensor products, equivariant direct summands,
and finite equivariant filtrations whose successive quotients belong to
the class.  An extension in this definition is an exact sequence of
compact semilinear $A_m$-modules.
\end{definition}

\begin{lemma}
\label[lemma]{lem:mark-powered-division}
The algebra $D_m(\ell,b)$ is finite free of rank $p^{n\ell-b}$ over
$A_m$.  Its construction and action are intrinsic to relative
Frobenius and commute with restriction to a deeper height locus.
Every member of $\mathscr C_m$ is finite free over $A_m$, and reduction
modulo $x$ sends $\mathscr C_m$ into $\mathscr C_{m+1}$.
\end{lemma}

\begin{proof}
Write $[p^\ell]_{\cF}(T)=f(T^{p^b})$.  Modulo the maximal ideal of
$A_m$, the series $f(Z)$ is $Z^{p^{n\ell-b}}$.  Weierstrass preparation
therefore gives
\[
 f(Z)=v(Z)W(Z),\qquad
 W(Z)=Z^{p^{n\ell-b}}+\sum_{i<p^{n\ell-b}}c_iZ^i,
\]
where $v$ is a unit and all $c_i$ lie in the maximal ideal.  Consequently
$[p^\ell](T)$ generates the same ideal as $W(T^{p^b})$.
Division by this monic polynomial shows both that
$A_m[Z]/(W(Z))\longrightarrow A_m[[T]]/([p^\ell](T))$, $Z\mapsto T^{p^b}$,
is injective and that its image has basis
$1,T^{p^b},\ldots,T^{p^b(p^{n\ell-b}-1)}$.  This proves the rank assertion.

A coordinate change $T\mapsto a(T)$ carries $T^{p^b}$ to
$a(T)^{p^b}$, a series in $T^{p^b}$.  Thus the subalgebra is preserved;
it is the relative Frobenius image of the division scheme on coordinate
rings.  The same monic-division calculation after base change identifies
the resulting algebra with the corresponding powered division algebra
of the restricted group.  In particular it commutes with passage to
$A_{m+1}=A_m/(x)$.  The numerical restriction on $b$ persists, since
$b\leq m\ell\leq(m+1)\ell$.

The basic modules are finite free.  An extension with finite free
quotient splits as an underlying $A_m$-module, so its middle term is
finite free.  A finite direct summand is projective and hence free over
the local ring $A_m$.  These facts also show that reduction modulo $x$
preserves the exact filtrations in the definition.  The periodicity line
restricts to the corresponding line at height $m+1$.  The last assertion
now follows by induction over the finite constructions defining the
class.
\end{proof}

\subsection{Connected jets and whole finite models}

For the remainder of this section assume $m<n$.
Over $B_m$, the connected part of $\cF[p^\infty]$ is a formal group of
height $m$.  Its isomorphisms to $\Gamma_m$ form a pro-finite-\'{e}tale
$P_m$-torsor.  The ring-level assertion used here is
\citet[Lecture~14, Theorem~1]{lurieFormalGroups}: the isomorphism algebra
of two formal group laws of the same exact positive height over a ring
is an increasing union of finite \'{e}tale algebras.  In the present
Honda trivialization its finite quotients can be chosen as follows.

\begin{proposition}[Integral connected-jet models]
\label[proposition]{prop:mark-jets}
For each $h\geq0$, the connected marking through exponent $p^h$ has a
finite free model $H_{m,h}$ over $A_m$.  There are coordinates
$a_0,\ldots,a_h$ and polynomials in earlier coordinates such that
\begin{equation}
\begin{aligned}
 a_0^{p^m-1}&=x,\\
 a_i^{p^m}-x^{p^i}a_i&=P_i(a_0,\ldots,a_{i-1}),
       &&1\leq i\leq h.
\end{aligned}
 \label{eq:mark-jet-equations}
\end{equation}
If the first Hasse coefficient is normalized only up to a unit, the
corresponding unit factors are inserted in these equations.  The
residue monomials
\[
 a_0^{e_0}a_1^{e_1}\cdots a_h^{e_h},\qquad
 0\leq e_0<p^m-1,\quad 0\leq e_i<p^m\ (i>0),
\]
are an $A_m$-basis.  The $G$-action preserves this model and admits a
finite filtration with quotients integer powers of $\cL_m$.  The
filtration remains such a filtration on every deeper height locus.
\end{proposition}

\begin{proof}
Write an isomorphism $f:\cF\longrightarrow\Gamma_m$ as an ordinary
power series.  Its coefficient at $T^{p^i}$ is $a_i$.  At a degree $j$
which is not a power of $p$, some intermediate binomial coefficient
$\binom{j}{v}$ is nonzero modulo $p$.  In the coefficient of
$X^vY^{j-v}$ in
$f(\cF(X,Y))=\Gamma_m(f(X),f(Y))$, the unknown coefficient of $T^j$
therefore has a nonzero constant coefficient.  All other terms involve
earlier coefficients.  It follows inductively that every non-$p$-power
coefficient is a polynomial over $A_m$ in the preceding $a_i$.

Now compare $f([p]_{\cF}(T))$ with $f(T)^{p^m}$.  In degree $p^m$ this
gives $a_0x=a_0^{p^m}$, hence the first equation, since $a_0$ is
invertible on the open marking space.  In degree $p^{m+i}$, the only
term involving $a_i$ on the left is $x^{p^i}a_i$; all remaining terms
involve earlier coefficients.  This proves
\cref{eq:mark-jet-equations}, with polynomials over $A_m$.  The
recursive comparison is also the usual proof of the Honda
classification; see \citet[Theorem~A2.2.12, its proof, and
Lemmas~A2.2.20--A2.2.21]{ravenel1986}.

Let $H_{m,h}$ be the successive monic algebra defined by these
equations.  Monic division gives the displayed basis, so it is finite
free of rank $(p^m-1)p^{mh}$.  After inverting $x$, $a_0$ is invertible,
the derivative of its equation is invertible, and the derivative of
the $i$th subsequent equation is $-x^{p^i}$.  Thus
$H_{m,h}[1/x]$ is finite \'{e}tale.

The actual marking quotient is finite \'{e}tale by the ring-level
classification above.  After a strictly henselian base change its
fiber is the corresponding quotient of $P_m$.  A Honda automorphism
has $p^m-1$ possible nonzero linear coefficients, and $p^m$ choices at
each succeeding $p$-power; the other coefficients are then determined.
This follows either from the Honda coefficient calculation or from
the order filtration on its endomorphism ring.  Hence that quotient
has rank $(p^m-1)p^{mh}$. The equations define a map from
$H_{m,h}[1/x]$ to its coordinate ring. Every jet coefficient is a
polynomial in $a_0,\ldots,a_h$, so these coordinates distinguish
geometric jets. The map on geometric function algebras is therefore
surjective, and finite \'{e}taleness gives surjectivity of the map
itself. A surjection between finite
\'{e}tale algebras of the same rank is an isomorphism, as can be checked
after strict henselization.  This identifies the algebra, including
over nonreduced test bases, with the connected marking quotient.

Precomposing $f$ with $\alpha_g$ gives
\begin{equation}
 g(a_i)=\chi_g^{p^i}a_i+Q_{g,i}(a_0,\ldots,a_{i-1}),
 \qquad g(a_0)=\chi_ga_0.
 \label{eq:mark-jet-triangular}
\end{equation}
Every coefficient in this formula lies in $A_m$; only finitely many
coefficients of $\alpha_g$ are used at a fixed jet order.  The formulas
preserve the defining relations after inverting $x$.  Since
$H_{m,h}$ is $x$-torsion-free, they preserve them before localization
as well.  The inverse coordinate change supplies the inverse action.

Order the displayed basis lexicographically starting with $e_h$, then
$e_{h-1}$, and so on.  In expanding a transformed monomial, any use of
an earlier-variable term in \cref{eq:mark-jet-triangular} lowers this
order.  Reducing a power $a_i^{p^m}$ by its monic equation lowers it
again.  The diagonal term on a basis vector is therefore
$\chi_g^{e_0+pe_1+\cdots+p^he_h}$.  The initial spans in this order are
stable, with the claimed line quotients.  All these spans and
quotients are free over $A_m$, so the filtration specializes exactly
to every deeper height locus.  Finally the formulas are continuous
in $g$ and in the compact coefficient topology, because every fixed
finite coefficient quotient uses finitely many coefficients of the
continuous universal coordinate change.
\end{proof}

On the same open, marking the \'{e}tale Tate module gives a
$\GL_r(\Zp)$-torsor.  Its action commutes with changes of connected
marking and with $G$.  Put
\[
 J=P_m\times\GL_r(\Zp).
\]
We use a cofinal sequence of normal product congruence subgroups
$U\subset J$.  The connected factors can be chosen among the kernels
of the jet quotients just constructed: the subsequence corresponding
to $1+p^a\cO_{D_m}$ is cofinal.  After omitting finitely many terms,
both factors, and hence $U$, are uniform and torsion-free.  Let $B_U$
denote the finite \'{e}tale $J/U$-torsor algebra of the two markings.

\begin{proposition}[A whole model for a finite marking]
\label[proposition]{prop:mark-whole-model}
For each such $U$ there are a finite free $A_m$-algebra $A'$, with
integral multiplication laws in a fixed basis, and a $G$-invariant
idempotent $e\in B'=A'[1/x]$ such that
\[
 B_U=eB',\qquad B'/B_m\text{ is finite \'{e}tale}.
\]
As a semilinear $A_m$-module, $A'$ has a finite filtration whose
quotients are basic coefficients in $\mathscr C_m$.  Its reduction
modulo $x$ has such a filtration in $\mathscr C_{m+1}$.
In a whole basis $A'=\bigoplus_\nu A_m b_\nu$, the $G$-action is
given by continuous matrices over $A_m$.  The base action is a
continuous formal action preserving $(x)$; in particular the
whole-basis action has a uniform pole bound, including for compact
families of group elements.
\end{proposition}

\begin{proof}
Choose the connected jet order $h$ and the \'{e}tale level $\ell$
defining $U$.  Relative Frobenius of exponent $m\ell$ kills precisely
the connected part of $\cF[p^\ell]$ on the exact-height open.  Its
image is the \'{e}tale quotient, whose coordinate algebra is
$D_m(\ell,m\ell)[1/x]$.  One can check \'{e}taleness directly: writing
$[p^\ell](T)$ in the variable $Z=T^{p^{m\ell}}$, its linear coefficient
is invertible when $x$ is invertible.  Equivalently its geometric
fibers have $p^{r\ell}$ distinct points, and the resulting finite
flat group scheme is \'{e}tale.

The set of ordered bases in the $r$-fold power of this finite
\'{e}tale group is an open and closed subscheme.  Indeed it can be
checked after an \'{e}tale trivialization of the group, where it is a
subset of a finite constant set.  Its characteristic idempotent is
preserved by every isomorphism of the group, in particular by $G$.
Take
\[
 A'=H_{m,h}\otimes_{A_m}D_m(\ell,m\ell)^{\otimes r}
\]
and extend that idempotent over the connected factor.  This gives
the asserted $e$ and the identification with the two marking
torsors.  Both factors of $B'$ are finite \'{e}tale over $B_m$.

Tensor the monomial filtration of \cref{prop:mark-jets} with the whole
powered tuple algebra.  Its quotients are precisely of the form
\cref{eq:mark-basic-coefficients}.  The same holds after reduction
modulo $x$ by \cref{lem:mark-powered-division}.  Tensoring the finite
free bases gives a basis of $A'$, and all multiplication constants
lie in $A_m$.  This construction uses the entire tuple algebra at
the boundary; no extension or specialization of $e$ across $x=0$
is involved.

The base action and the action on a powered division coordinate
are substitution by the universal formal coordinate change and
its appropriate Frobenius power.  Their reduction in a fixed
finite free basis converges in the parameter topology.  Modulo
any power of the maximal ideal, only a finite truncation is
needed, so its matrix entries vary continuously in $g$.  The
connected-jet entries have the same property by
\cref{eq:mark-jet-triangular}.  Thus all matrices have integral
entries and zero pole loss in the whole lattice, uniformly in
$g$.  Preservation of $(x)$ is \cref{eq:mark-height-character}.
\end{proof}

\subsection{Generator lifts and roots}

At full markings the schemes of lifts of the ordered \'{e}tale
generators are torsors under the connected division groups.  Their
inverse limit is a torsor for
\[
 T=\varprojlim_{[p]}\Gamma_m[p^\ell]^r.
\]
Here $T$ and its finite quotients are affine group schemes; their
representations below are coactions of the coordinate Hopf algebras.
The two marking groups act on this description by changing the two
bases; $G$ commutes with the resulting translations.

\begin{proposition}[The lift algebras are root algebras]
\label[proposition]{prop:mark-roots}
At level $\ell$, the reduced closure over $A_m$ of lifted geometric
\'{e}tale bases has coordinate algebra
\[
 R_{m,\ell}=\kk[[t_{1,\ell},\ldots,t_{r,\ell}]].
\]
If $s_i=t_{i,\ell}^{p^{m\ell}}$, the intermediate algebra generated
by the $s_i$ over $A_m$ is
\[
 S_{m,\ell}=\kk[[s_1,\ldots,s_r]],\qquad
 A_m\subseteq S_{m,\ell}.
\]
After inverting $x$, $S_{m,\ell}[1/x]$ is the \'{e}tale basis algebra,
and $R_{m,\ell}[1/x]$ is its full generator-lift algebra.  In
particular the latter is the root extension of order $p^{m\ell}$,
finite free of rank $p^{m\ell r}$, with no reduction of the lift
scheme omitted.  These identifications are natural under the $G$
action, changes of markings, and finite \'{e}tale base change.
For $j\leq\ell$, earlier division coordinates are integral series
in the $p^{m(\ell-j)}$ powers of the level-$\ell$ coordinates.
In particular the base parameters, in first-level coordinates,
have order at least $p^m$.
\end{proposition}

\begin{proof}
Strauch describes the corresponding torsion stratum over the completed
maximal unramified coefficient base.  There the full Drinfeld
level-$p^\ell$ deformation ring $\widetilde R_\ell$ is finite flat
over the Lubin--Tate deformation ring and regular, with the universal
torsion coordinates $\phi_1,\ldots,\phi_n$ as regular parameters.
The quotient
\[
 \widetilde R_\ell/(\phi_1,\ldots,\phi_m)
\]
is a regular ring of dimension $r$ over the height-at-least-$m$
base; its generic torsion coordinates $\phi_{m+1},\ldots,\phi_n$
give the lifted \'{e}tale bases.  The induced generic Galois action
is the full $\GL_r(\Z/p^\ell)$ action
\citep[\S3, Proposition~4.1, and the proof of Theorem~2.1]{strauch2007}.
Over our finite Honda field we establish the whole reduced closure
directly, using the division divisor.

Let $C_\ell$ be the coordinate algebra of the whole $r$-fold division
tuple scheme over $A_m$.  Over $B_m$, let $Q_\ell$ be the finite
\'{e}tale algebra of ordered bases of the \'{e}tale division group,
and let $L_\ell$ be its full generator-lift algebra.  The map
$Q_\ell\to L_\ell$ is finite faithfully flat: its fibers are torsors
under the product of the $r$ connected division groups.  The basis
condition selects an open and closed subscheme of the powered tuple
scheme, so $C_\ell[1/x]\to L_\ell$ is surjective.  Define
\[
 R=\operatorname{image}\bigl(C_\ell\longrightarrow (L_\ell)_{\mathrm{red}}\bigr).
\]
This is the reduced closure of the entire lifted basis locus.  It is
finite over $A_m$, as a quotient of $C_\ell$, and
$R[1/x]=(L_\ell)_{\mathrm{red}}$.  Every component dominates $B_m$:
the finite flat algebra $L_\ell/B_m$ has its minimal primes over the
generic point.  Taking their closures adds no vertical component.
Consequently $A_m\to R$ is injective and integrality gives $\dim R=r$.
The closed fiber of $C_\ell$ is supported at the tuple of origins, so
$R$ is complete local with residue field $\kk$.
Denote the universal lifted
coordinates by $t_i=t_{i,\ell}$, and put
$t_{i,1}=[p^{\ell-1}]_{\cF}(t_i)$.  At geometric generic points, the
$p^r$ combinations of the $t_{i,1}$ are all the points of
$\cF[p]$, each with multiplicity $p^m$.  Weierstrass preparation
therefore gives the divisor identity
\begin{equation}
 [p]_{\cF}(Z)
  =v(Z)\prod_{a\in\Fp^r}
       \left(Z-\sum_{i=1}^{r}{}_{\cF}[a_i](t_{i,1})\right)^{p^m}
       \quad\text{in }R[[Z]],
 \label{eq:mark-basis-divisor}
\end{equation}
where $v$ is a unit.  Equality holds in $R$, since $R$ is reduced and
all its components meet the generic basis locus.  Modulo $(t_1,\ldots,t_r)$,
the right side is a unit times $Z^{p^n}$.  Successively comparing the
Hasse coefficients forces every $u_m,\ldots,u_{n-1}$ to vanish in
that quotient.  Consequently the maximal ideal of $R$ is generated
by the $r$ elements $t_i$.  Completeness gives a surjection
$\kk[[T_1,\ldots,T_r]]\twoheadrightarrow R$.  A nonzero ideal of
the power-series domain has quotient of dimension at most $r-1$;
since $\dim R=r$, the kernel is zero.  This proves the asserted
description of $R_{m,\ell}$ over $\kk$.

For the powered assertion, set $h=p^{m\ell}$ and let
$S=A_m[s_1,\ldots,s_r]\subset R$, where $s_i=t_i^h$.  This is a
finite $A_m$-module: it is a submodule of a finite module over the
Noetherian ring $A_m$.  It is therefore complete; it is local with
residue field $\kk$, and $\dim S=r$ by integrality.  The first-level
coordinates raised to $p^m$ belong to $S$.  Indeed,
$[p^{\ell-1}](T)$ is a series in $T^{p^{m(\ell-1)}}$, so raising
$[p^{\ell-1}](t_i)$ to $p^m$ gives an integral series in $s_i$.
The same is true of all their formal-group combinations.

To apply \cref{eq:mark-basis-divisor} inside $S$, write
$[p]_{\cF}(Z)=V(Z^{p^m})$.  The monic factor on its right is a
polynomial in $Z^{p^m}$ whose coefficients lie in $S$.  Division of
$V(W)$ by this monic polynomial in $W$ is performed in $S[[W]]$;
the remainder is zero after embedding in $R[[W]]$, hence is already
zero in $S[[W]]$.  The quotient is a unit, as its reduction at the
closed point is a unit.  Reducing now modulo $(s_1,\ldots,s_r)$
again gives a unit times $W^{p^r}$.  The successive Hasse
coefficients therefore put all the base parameters in $(s_1,\ldots,s_r)S$.
Thus $S$ is complete with maximal ideal generated by the $r$
elements $s_i$.  The dimension argument just used gives
$S=\kk[[s_1,\ldots,s_r]]$.  In particular the inclusion of the base
factors through this power-series subring.

The powered coordinates on the exact-height open are the \'{e}tale
division coordinates.  Their basis condition gives a surjection
$Q_\ell\to S[1/x]$, since the $s_i$ generate the target over $B_m$.
The composite into $R[1/x]=(L_\ell)_{\mathrm{red}}$ is injective:
$Q_\ell$ is reduced, and the faithfully flat map
$Q_\ell\to L_\ell$ cannot send a nonzero element of $Q_\ell$ to
a nilpotent.  Thus $Q_\ell\simeq S[1/x]$, identifying the basis
algebra on the whole open.

The extension
\[
 \kk[[s_1,\ldots,s_r]]\ \longrightarrow\
 \kk[[t_1,\ldots,t_r]],\qquad s_i\longmapsto t_i^h,
\]
is free with basis $\prod_i t_i^{e_i}$, $0\leq e_i<h$; this follows
by partitioning the monomials in a power series by their exponents
modulo $h$.  Its rank is $h^r$.  After localization, $R[1/x]$ is a
closed subscheme of the generator-lift scheme over $S[1/x]$.  That
lift scheme is a torsor under the product of $r$ connected
$p^\ell$-division groups, so it too is finite locally free of rank
$p^{m\ell r}=h^r$.  On coordinate rings this is a surjection of
finite locally free modules of equal rank, hence an isomorphism.
This proves the assertion about the full scheme, including its
nonreduced geometric fibers.

All maps used to define $s_i$ are relative Frobenius maps, and all
transition maps are multiplication by powers of $p$.  They commute
with group isomorphisms.  Changing a basis replaces the $t_i$ by
formal integral linear combinations, whose $h$th powers are the
corresponding combinations of the $s_i$.  Hence the descriptions
are natural for $G$ and both marking actions.  For an \'{e}tale
algebra, relative Frobenius is an isomorphism
\cite[Lemma~41.14.3, Tag~0EBS]{stacks}.  Root extension therefore
commutes with any subsequent finite \'{e}tale marking extension:
if $E/S[1/x]$ is \'{e}tale, then
\[
 E\otimes_{S[1/x]}R[1/x]\simeq E^{1/h}.
\]
The same identity gives descent through marking levels, with their
natural change-of-marking actions retained.  Finally
$t_{i,j}=[p^{\ell-j}](t_{i,\ell})$ is an integral series in
$t_{i,\ell}^{p^{m(\ell-j)}}$ with coefficients in $A_m$.
Substituting the already proved expressions for those coefficients
puts it in
$\kk[[t_{1,\ell}^{p^{m(\ell-j)}},\ldots,
 t_{r,\ell}^{p^{m(\ell-j)}}]]$.
At level one, each base parameter lies in the maximal ideal of
$\kk[[t_{1,1}^{p^m},\ldots,t_{r,1}^{p^m}]]$, which gives the stated
order bound.
\end{proof}

\begin{lemma}[Basic coefficients on the marked torsor]
\label{lem:mark-associated-coefficients}
At full markings, each basic coefficient in
\cref{eq:mark-basic-coefficients} is associated to a finite-dimensional
translation representation over $\kk$.  Its torsor, representation,
and maps descend to finite marking levels, with the natural
change-of-marking actions.  A fixed finite collection of such data
admits a common sufficiently deep level.  The periodicity line
already trivializes at the linear connected jet.
\end{lemma}

\begin{proof}
Choose a level $L$ at least as large as all the $\ell_j$.  A lift
of the marked \'{e}tale generators in $\cF[p^L]$ supplies a splitting
of the connected--\'{e}tale extension at that level: send a standard
basis vector of $(\Z/p^L)^r$ to its lift and extend by the group
law.  The lifts are $p^L$-torsion, so this is well defined.  The
resulting product decomposition identifies the division group with
\[
 \Gamma_m[p^L]\times(\Z/p^L)^r.
\]
Changing the lifts translates this splitting by
$\Gamma_m[p^L]^r$.  The coordinate algebra of the displayed group
is a finite-dimensional $\kk$-space; its restrictions to the lower
levels and their powered-coordinate images are functorial under
relative Frobenius.  They therefore give finite-dimensional
representations of the same finite translation quotient.  Torsor
descent identifies their associated bundles with the original
powered division coefficients.  Tensor products and finite sums
preserve this identification.

The derivative of the connected marking identifies the intrinsic
linear coefficient line with the Honda line.  This is an
equivariant trivialization and uses only $a_0$, which is invertible
on the open.  It treats every positive or negative periodicity
power.  The schemes and their coactions use finite division level,
finitely many coefficients, and a finite-dimensional vector space
over the finite field $\kk$.  The continuous change-of-marking
action consequently factors through a finite quotient on each
fixed datum.  Finite presentation of the torsor descent and its
maps puts them at a finite marking stage; the maximum of finitely
many stages works for a finite collection.  This assertion is about
the basic coefficients.  Arbitrary extensions and summands in
$\mathscr C_m$ will be handled by exact sequences and retracts,
without asserting that every such extension is itself a constant
translation representation.
\end{proof}

\subsection{Finite isotropy and trace one}

\begin{proposition}[Geometric hypotheses for the isotropy bound]
\label[proposition]{prop:mark-finite-isotropy}
Suppose $U$ is a torsion-free product marking subgroup as above.
Every finite subgroup $F\subset G$ acts geometrically freely on
$\Spec B_U$.  There is an element $a_F\in B_U$ with
\[
 \sum_{f\in F}f(a_F)=1.
\]
Multiplication by $a_F$ is continuous between bounded-denominator
steps of finite coefficient bundles.  It also acts continuously
on bounded analytic section and compact-support models, with
compact profinite parameters retained.  Thus the geometric
coefficient modules satisfy the trace and module-action
hypotheses of \cref{prop:cts-isotropy-bound,prop:cts-full-duality}.
Equivalently they satisfy \cref{def:cts-trace}.
\end{proposition}

\begin{proof}
Let $g\ne1$ have finite order $N$, and suppose it fixes a geometric
point $z$ of $\Spec B_U$.  Its image on the base factors through
\[
 R=A_m/I_g,\qquad I_g=(g(a)-a:a\in A_m).
\]
Every ideal of a Noetherian complete local ring is closed, so $R$
is complete local.  The framing action becomes an actual
automorphism $\gamma$ of the formal group over $R$.  Its reduction
at the closed point is the nonidentity Honda automorphism $g$.
The difference endomorphism
\[
 f(T)=\gamma(T)-_{\cF}T
\]
therefore has nonzero reduction, of some finite order $d$ in $T$.
Weierstrass preparation over $R$ makes
$H=\Spec R[[T]]/(f(T))$ finite free of rank $d$.
For every $\ell$ the fixed subscheme of the finite division scheme is
\[
 H_\ell=\Spec R[[T]]/(f(T),[p^\ell](T)),
\]
a closed subscheme of $H$.  Every geometric fiber of $H_\ell$ thus
has coordinate dimension at most $d$.  This preparation and these
finite quotients were formed over the complete base before
localization.

Lift $z$ to a full marking $t$ over its algebraically closed residue
field.  Such a lift exists by the marking torsors: the inverse
system of nonempty finite fibers has surjective transitions.  As
$g$ fixes the $U$-coset, $gt=tu$ for some $u\in U$.  Commutation
of the actions gives $u^N=1$.  Since $U$ is torsion-free, $u=1$.
In particular $\gamma$ induces the identity on the connected
height-$m$ formal group, hence on its finite $p^\ell$-division
subgroup inside the specialized division scheme.  That subgroup
has rank $p^{m\ell}$ and is contained in $(H_\ell)_z$.  Taking
$p^{m\ell}>d$ is a contradiction.  Notice that only the connected
marking is needed for this last contradiction; no splitting of
the connected--\'{e}tale extension at $z$ is being used.
To spell out the specialization, its connected division coordinate
is nilpotent, so evaluation of each original formal series there
is a finite sum.  It gives a surjection from the specialized
prepared algebra defining $H_\ell$ onto that connected coordinate
algebra, since both $f$ and $[p^\ell]$ vanish there.  This uses only
the finite algebra formed over $R$ and never identifies
$R[[T]]\otimes_R\Omega$ with $\Omega[[T]]$.

We give the trace argument for an arbitrary commutative algebra
$D$ with a geometrically free action of a finite group $F$; this
avoids a finite-type assumption on $B_U$.  Put $C=D^F$.  The image
$I=\Tr_F(D)$ is an ideal of $C$, because trace is $C$-linear.  If
$I$ were contained in a maximal ideal $\mathfrak m$ of $C$, the
integrality of $D/C$ would give a point $z:D\to\Omega$ extending
$C\to\kappa(\mathfrak m)$, for an algebraically closed field
$\Omega$.  Integrality follows directly from the orbit polynomials
$\prod_{f\in F}(T-f(d))$, and the point follows by lying-over.
In $D_\Omega=D\otimes_C\Omega$, consider the surjective
$\Omega$-algebra evaluations
\[
 e_f(d\otimes\lambda)=z(f(d))\lambda,\qquad f\in F.
\]
Geometric freeness says that these maps are distinct.  Their
maximal kernels are therefore distinct, and the Chinese remainder
theorem gives $b\in D_\Omega$ such that
$e_1(b)=1$ and $e_f(b)=0$ for $f\ne1$.  Hence
$e_1(\Tr_F(b))=1$.  On the other hand, writing
$b=\sum_i d_i\otimes\lambda_i$ gives
\[
 \Tr_F(b)=\sum_i\Tr_F(d_i)\otimes\lambda_i=0,
\]
because all $\Tr_F(d_i)$ lie in $\mathfrak m$.  This contradiction
shows that $I=C$.  Thus a single element of $D$ has trace one.
The argument uses no assertion that invariants commute with
residue-field base change.  Apply it with $D=B_U$.

Finally, write $a_F=e x^{-c}a$ in the whole finite model, with
$a\in A'$ and $c$ finite.  Multiplication is a finite matrix over
$B_m$ in the whole basis.  Clearing its finitely many denominators
makes it a continuous map from any compact lattice into a fixed
enlarged lattice.  The same statement holds for finite coefficient
bundles, after choosing lattices.  On an analytic bound where
$x$ is bounded away from zero, the same function and finite
matrices are bounded.  Multiplication is continuous on sections,
preserves supports, and commutes with restriction; it therefore
acts on the section and support complexes.  Retaining a compact
profinite parameter means using the uniform seminorm on that
parameter, so the same bound applies.  These are precisely the
trace-one multiplication and semilinearity conditions in the
cited cohomological propositions.  The additional continuity of
the later Laurent representatives under changes of coordinates
is verified in \cref{lem:laurent-action}.
\end{proof}

\subsection{Boundary strips}

\begin{lemma}[Boundary comparison under the next-height hypothesis]
\label[lemma]{lem:mark-strips}
Assume that $H^i_\cts(G,M)$ is finite for every $M\in\mathscr C_{m+1}$
and every $i\geq0$.  Let $L$ be either a member of $\mathscr C_m$ or
$A'\otimes_{A_m}M$ for $M\in\mathscr C_m$, where $A'$ is a whole
model from \cref{prop:mark-whole-model}.  For any integers $a<b$,
\[
 H^i_\cts(G,x^aL/x^bL)
\]
is finite in every degree.  For each fixed integer $a$, localization
induces a map
\[
 H^i_\cts(G,x^aL)\longrightarrow H^i_\cts(G,L[1/x])
\]
with countable kernel and cokernel.  These assertions concern
whole lattices; they do not require a model for the idempotent
$e$ on the closed locus.
\end{lemma}

\begin{proof}
The quotient $x^aL/x^bL$ has a finite filtration whose layers are
$x^jL/x^{j+1}L$ for $a\leq j<b$.  Multiplication by $x^j$ and
\cref{eq:mark-height-character} identify the $j$th layer with
\[
 (L/xL)\otimes_{A_{m+1}}
       \cL_{m+1}^{\otimes j(p^m-1)}.
\]
If $L\in\mathscr C_m$, this belongs to $\mathscr C_{m+1}$ by
\cref{lem:mark-powered-division}.  For a whole marking model it
has a finite filtration with the same property by
\cref{prop:mark-whole-model}.  The next-height hypothesis and the
long exact cohomology sequences of these finite filtrations prove
finiteness of every strip.

Put $V=L[1/x]/x^aL$.  It is the countable increasing union of the
finite-width strips $x^{a-j}L/x^aL$, $j\geq0$.  With the stipulated
bounded-denominator cochains, its cochain complex is their
filtered colimit.  Filtered colimits of vector spaces are exact,
so its cohomology is the filtered colimit of their finite
cohomology groups.  It is countable, since $\kk$ is finite.  The
short exact coefficient sequence
\[
 0\longrightarrow x^aL\longrightarrow L[1/x]
   \longrightarrow V\longrightarrow0
\]
is exact on continuous cochains by \cref{lem:cts-exact}; in
concrete terms one lifts a cochain through a compact strip after
enlarging its denominator bound.  The long exact sequence
identifies the localization kernel with an image of
$H^{i-1}_\cts(G,V)$ and its cokernel with a subspace of
$H^i_\cts(G,V)$.  Both are countable, as claimed.
\end{proof}

\paragraph{The remaining induction step.}
At $m=n$, the allowed coefficients are finite-dimensional over $\kk$,
so the finite-type resolution of \cref{lem:cts-comparison} gives finite
cohomology in each degree. For $m<n$, assume this finiteness for every
coefficient in $\mathscr C_{m+1}$. The target of
\crefrange{sec:completed-tower}{sec:pro-transfer}
is the marked scalar assertion
\[
 H^i_{\cts}(G,B_{U_\nu})\text{ is countable}
 \qquad(i\geq0,\ \nu\geq\nu_0),
\]
where $(U_\nu)$ is a cofinal decreasing sequence of normal product
congruence subgroups and one $\nu_0$ works for all degrees.
This is \cref{thm:pro-countability}: it concerns each individual
finite marking level on that tail.

The return to arbitrary allowed coefficients still requires
\cref{sec:coefficients}. There the associated translation
representations extend scalar countability to basic coefficients,
and finite marking descent removes the cover. The boundary comparison
above then makes the original compact coefficient cohomology
countable. It is compact Hausdorff by \cref{lem:cts-comparison}, hence
finite. Exact filtrations and summands give the full allowed class,
closing the induction in \cref{thm:coeff-finite}.
